\begin{tcolorbox}[enhanced,height=7.1cm,colback=navy,colframe=navy,boxrule=0pt,
  arc=0pt,valign=center,left=14mm,right=14mm]
\color{white}
\begin{minipage}[c]{.27\textwidth}\centering
{\fontsize{76}{80}\selectfont\bfseries 02}\par\vspace{2mm}
{\color{cyanlearn}\rule{.82\linewidth}{1.5pt}}
\end{minipage}\hfill
\begin{minipage}[c]{.66\textwidth}
{\fontsize{34}{39}\selectfont\bfseries Learning Journey \&\par Courses Plan}
\end{minipage}
\end{tcolorbox}

\begin{center}
\textcolor{red}{\bfseries\itshape STRICTLY CONFIDENTIAL}\qquad
\textbf{Prepared by:} SBM ITB --- Direktorat Pendidikan Profesional Berkelanjutan\par
Kementerian Pertahanan Republik Indonesia \textbullet{} DIKTISAINTEK BERDAMPAK \textbullet{} Danantara Indonesia
\end{center}

\section{Learning Journey: 5 Components for Transformation}
\textit{Ekosistem 5-komponen terintegrasi untuk mengubah potensi pelajar menjadi
\textbf{strategic business leader} dengan kemampuan memimpin dan eksekusi.}

\begin{tcolorbox}[colback=lightblue,colframe=navy,boxrule=.8pt,arc=2mm,
  title={Integrated Design Logic},colbacktitle=navy,coltitle=white,fonttitle=\bfseries]
Program ini tidak hanya terdiri dari kelas \textit{stand-alone}, namun mengintegrasikan kurikulum, metode pembelajaran, projek bisnis real, mentoring, teknologi, dan lingkungan kampus residensial dalam satu arsitektur pengembangan.
\end{tcolorbox}

\begin{infobox}[1. Desain Kurikulum --- \textit{3 Learning Building Blocks}]{brick}
\begin{tcbraster}[raster columns=3,raster equal height=rows,raster column skip=2mm]
\begin{tcolorbox}[colback=palered,colframe=palered,boxrule=0pt]
\textbf{1 -- Foundation Learning}\par
Learning agility \textbullet{} Communication \textbullet{} Professional presence \textbullet{} Strategic thinking \textbullet{} Basic accounting \textbullet{} Management fundamentals
\end{tcolorbox}
\begin{tcolorbox}[colback=palered,colframe=palered,boxrule=0pt]
\textbf{2 -- 6 Core Courses}\par
6 MK inti: OB \& MP, Financial Management \& Strategy, Marketing Management, OSCM, Strategic Decision Making \& Negotiation dalam 6 sektor BUMN
\end{tcolorbox}
\begin{tcolorbox}[colback=palered,colframe=palered,boxrule=0pt]
\textbf{3 -- Managerial, Leadership, Entrepreneurship}\par
Managerial skill \textbullet{} Leadership skill \textbullet{} Entrepreneurship as capability to lead, manage, and create value
\end{tcolorbox}
\end{tcbraster}
\end{infobox}

\begin{infobox}[2. Metode Pembelajaran --- \textit{Case-based \& blended exposure}]{red}
\begin{tcbraster}[raster columns=3,raster equal height=rows,raster column skip=2mm]
\begin{tcolorbox}[colback=palered,colframe=palered,boxrule=0pt]
\textbf{Core Method}\par Pembelajaran case-based \textbullet{} Studi kasus \textbullet{} Diskusi \textbullet{} Simulasi \textbullet{} Gamifikasi \textbullet{} Refleksi
\end{tcolorbox}
\begin{tcolorbox}[colback=palered,colframe=palered,boxrule=0pt]
\textbf{Learning Reinforcement}\par Pra-pembelajaran via LMS \textbullet{} Rencana studi dan rekap mingguan \textbullet{} Leader talks \textbullet{} Business sharing \textbullet{} Company visit
\end{tcolorbox}
\begin{tcolorbox}[colback=palered,colframe=palered,boxrule=0pt]
\textbf{Perspektif}\par Kombinasi academic rigor + business relevance + national context + global perspective
\end{tcolorbox}
\end{tcbraster}
\end{infobox}

\begin{infobox}[3. Action Learning Project --- \textit{Isu bisnis nyata BUMN}]{orange}
\begin{center}
\begin{tikzpicture}[stage/.style={draw=orange!35,fill=paleorange,rounded corners=2mm,
 minimum height=12mm,text width=3.2cm,align=center,font=\bfseries},
 smart/.style={draw=orange!35,fill=paleorange,rounded corners=2mm,minimum height=12mm,
 text width=7.2cm,align=left},>={Latex[length=2.2mm]},node distance=3mm]
\node[stage] (s1) {Problem\\Scoping};
\node[stage,right=of s1] (s2) {Data\\Gathering};
\node[stage,right=of s2] (s3) {Analysis \&\\Solution Design};
\node[stage,right=of s3] (s4) {Coaching\\Review};
\node[stage,right=of s4] (s5) {Final\\Presentation};
\node[smart,right=of s5] (sm) {\textbf{SMART Gate \& Impact}\\
\textcircled{S}pecific \textcircled{M}easurable \textcircled{A}chievable
\textcircled{R}elevant \textcircled{T}ime-Bound\\[-1mm]
\textit{Output:} Rekomendasi implementasi, estimasi business impact, dan rencana follow-up};
\draw[->,thick,orange] (s1)--(s2); \draw[->,thick,orange] (s2)--(s3);
\draw[->,thick,orange] (s3)--(s4); \draw[->,thick,orange] (s4)--(s5);
\draw[->,thick,orange] (s5)--(sm);
\end{tikzpicture}
\end{center}
\end{infobox}

\begin{infobox}[4. Sistem Mentoring --- \textit{Arsitektur pendukung berlapis}]{gold}
\begin{tcbraster}[raster columns=5,raster equal height=rows,raster column skip=2mm]
\begin{tcolorbox}[colback=palegold,colframe=palegold,boxrule=0pt]\textbf{SBM ITB}\par Kedalaman akademik \& praktisi\end{tcolorbox}
\begin{tcolorbox}[colback=palegold,colframe=palegold,boxrule=0pt]\textbf{Mentor Bisnis BUMN}\par Konteks industri \& realitas eksekusi\end{tcolorbox}
\begin{tcolorbox}[colback=palegold,colframe=palegold,boxrule=0pt]\textbf{Business Coach}\par Kualitas proses berpikir \& perkembangan kepemimpinan\end{tcolorbox}
\begin{tcolorbox}[colback=palegold,colframe=palegold,boxrule=0pt]\textbf{Pengasuh}\par Menjaga disiplin, karakter, \& konsistensi berperilaku\end{tcolorbox}
\begin{tcolorbox}[colback=palegold,colframe=palegold,boxrule=0pt]\textbf{Project Owner}\par Memastikan isu yang dikerjakan relevan bagi BUMN\end{tcolorbox}
\end{tcbraster}
\end{infobox}

\begin{infobox}[5. Operasional \& Fasilitas --- \textit{Infrastruktur pembelajaran tepercaya}]{green}
\begin{tcbraster}[raster columns=2,raster equal height=rows,raster column skip=2mm]
\begin{tcolorbox}[colback=palegreen,colframe=palegreen,boxrule=0pt]
\textbf{Operasional Program}\par Jadwal \textbullet{} Alokasi kelas \textbullet{} Fakultas \& fasilitator \textbullet{} Mentoring/coach system \textbullet{} LMS/tools \textbullet{} Evaluasi \textbullet{} Koordinasi logistik
\end{tcolorbox}
\begin{tcolorbox}[colback=palegreen,colframe=palegreen,boxrule=0pt]
\textbf{Nawasena Living-Learning Environment}\par Jadwal \textbullet{} Auditorium \textbullet{} Perpustakaan \textbullet{} Klinik 24 jam \textbullet{} Ahli nutrisi \textbullet{} Konseling \textbullet{} Workshop \textbullet{} Akomodasi \textbullet{} Makan \textbullet{} Olahraga \textbullet{} Laundry \textbullet{} Asistensi partisipasi
\end{tcolorbox}
\end{tcbraster}
\end{infobox}

\begin{center}
\begin{tikzpicture}[input/.style={rounded corners=1.5mm,minimum width=3.8cm,
 minimum height=9mm,text=white,font=\bfseries,align=center},
 target/.style={rounded corners=1.5mm,draw=navy,fill=lightblue,minimum width=12.7cm,
 minimum height=14mm,text width=12cm,align=center,font=\bfseries},>={Latex[length=2.5mm]}]
\node[input,fill=brick] (a) at (0,0) {Desain Kurikulum};
\node[input,fill=red] (b) at (4.2,0) {Metode Pembelajaran};
\node[input,fill=orange] (c) at (8.4,0) {Action Learning Project};
\node[input,fill=gold] (d) at (12.6,0) {Sistem Mentoring};
\node[input,fill=green] (e) at (16.8,0) {Operasional \& Fasilitas};
\node[target] (t) at (8.4,-2) {\textcolor{navy}{Target Outcome:} Partisipan yang dapat memahami persoalan, mengambil keputusan yang bertanggung jawab, bekerja lintas fungsi, dan menghasilkan solusi yang berdampak\\[1mm]\textit{LEARN \textbullet{} LEAD \textbullet{} EXECUTE.}};
\foreach \x in {a,b,c,d,e}{\draw[->,thick,navy] (\x.south)--(\x.south |- t.north);}
\end{tikzpicture}
\end{center}

\begin{visualexplanation}
Lima komponen berwarna adalah masukan paralel menuju satu hasil. Desain kurikulum menetapkan fondasi, enam mata kuliah inti, dan kapabilitas managerial--leadership--entrepreneurship. Metode pembelajaran mengubah isi tersebut menjadi pengalaman case-based dan blended. Action Learning Project menjalankan urutan dari problem scoping sampai final presentation melalui gerbang SMART. Sistem mentoring memasangkan kedalaman akademik, konteks BUMN, coaching, pengasuhan, dan kepemilikan proyek. Operasional serta fasilitas menjaga semua aktivitas dapat berlangsung di lingkungan living-learning Nawasena. Seluruh panah bertemu pada target outcome yang sama: partisipan memahami persoalan, bertanggung jawab dalam keputusan, bekerja lintas fungsi, dan menghasilkan solusi berdampak melalui pola \textit{learn--lead--execute}.
\end{visualexplanation}

\section{Learning Journey: 70--20--10 Learning Model}
\subsection{Kegiatan Pembelajaran}
\Needspace{11cm}
\begin{tcbraster}[raster columns=4,raster equal height=rows,raster column skip=2mm,raster row skip=2mm]
\begin{infobox}[Weekly Study Plan/Recap]{red}
Peserta melakukan asesmen dan evaluasi terhadap materi pembelajaran pada minggu sebelumnya, dan melakukan rencana studi pada minggu tersebut.

\textit{Dilakukan di Kampus DCU Nawasena, didampingi oleh asisten kelas.}
\end{infobox}
\begin{infobox}[Pre-Learning]{red}
Peserta melakukan persiapan mandiri dengan mempelajari materi, studi kasus, maupun diskusi bersama grup untuk topik yang akan dibawakan pada minggu tersebut.

\textit{Dilakukan di Kampus DCU Nawasena, didampingi oleh asisten kelas.}
\end{infobox}
\begin{infobox}[Lecturing]{orange}
Tim pengajar (Dosen) dari SBM ITB membawakan materi dengan skema pembelajaran luring di setiap kelas terkait dengan topik sesuai dengan jadwal kurikulum.

\textit{Pengajaran secara luring oleh Dosen SBM ITB, didampingi oleh asisten kelas, didukung oleh asisten mata kuliah.}
\end{infobox}
\begin{infobox}[Case Study \& Roleplay]{gold}
Peserta akan diberikan \textit{case study paper} untuk dipelajari, didiskusikan bersama grup, dan dikerjakan untuk dilakukan pembahasan dan evaluasi kasus.

\textit{Didampingi oleh asisten kelas, modul menggunakan case study Harvard Business School.}
\end{infobox}
\begin{infobox}[Group Discussion ALP]{green}
Peserta melakukan diskusi di dalam grup untuk mempersiapkan kegiatan ALP terhadap kasus nyata dengan menerapkan teori yang sudah didapatkan.

\textit{Rangkaian ALP akan didampingi oleh asisten kelas, project owner, dan coach/mentor.}
\end{infobox}
\begin{infobox}[Group Discussion Social Project]{teal}
Peserta melakukan diskusi di dalam grup untuk mempersiapkan proyek sosial yang akan dilaksanakan pada bulan ke-4 untuk menerapkan teori yang sudah didapatkan.

\textit{Rangkaian SP akan didampingi oleh asisten kelas, pengasuh, dan coach/mentor.}
\end{infobox}
\begin{infobox}[Kegiatan Khusus]{blue}
Peserta mengikuti pembelajaran pelengkap dengan berbagai metode, seperti \textit{Executive Seminar}, \textit{Guest Lecture CEO BUMN}, \textit{Panel Discussion}, maupun \textit{Reflection Session}.

\textit{Kegiatan khusus dapat dilakukan dalam auditorium bersama seluruh peserta, didampingi oleh asisten kelas.}
\end{infobox}
\begin{infobox}[Business Immersion]{blue}
Peserta akan turun langsung keluar kelas untuk belajar secara imersif dengan sistem praktik/magang dan \textit{company visit} BUMN.

\textit{Kegiatan dilakukan di luar Kampus DCU Nawasena; akan dilakukan asesmen dengan laporan dan presentasi akhir.}
\end{infobox}
\end{tcbraster}

\begin{visualexplanation}
Delapan panel menunjukkan rentang pengalaman belajar dari persiapan individual sampai praktik di luar kampus. Weekly Study Plan/Recap dan Pre-Learning membingkai kegiatan mingguan; lecturing memberi fondasi formal; case study, roleplay, dan dua jenis group discussion memindahkan teori ke interaksi dan tindakan; kegiatan khusus mempertemukan seluruh cohort dengan praktisi; Business Immersion menempatkan peserta langsung dalam konteks BUMN. Teks miring pada setiap panel menjelaskan lokasi dan pihak pendamping yang menopang kegiatan tersebut.
\end{visualexplanation}

\subsection{Alokasi Sesi dan Kategori Metode Pembelajaran}
\begin{center}\small
\begin{tabularx}{.98\textwidth}{|L{5.5cm}|M{3.1cm}|M{2.2cm}|M{3.8cm}|Y|}\hline
\rowcolor{navy}\color{white}\bfseries Kegiatan & \color{white}\bfseries Jumlah Sesi/Minggu &
\color{white}\bfseries Total Minggu & \color{white}\bfseries Total Sesi Selama Program &
\color{white}\bfseries Kategori Metode Pembelajaran\\ \hline
Weekly Study Plan/Recap&2&14&28&\cellcolor{yellowlearn}Experiential Learning\\ \hline
\multirow{2}{*}{Pre-Learning}&2&14&28&\cellcolor{cyanlearn}Formal Learning\\ \cline{2-5}
&10&14&140&\cellcolor{greenlearn}Social Learning\\ \hline
Lecturing&6&14&84&\cellcolor{cyanlearn}Formal Learning\\ \hline
\multirow{2}{*}{Case Study}&2&14&28&\cellcolor{greenlearn}Social Learning\\ \cline{2-5}
&10&14&140&\cellcolor{yellowlearn}Experiential Learning\\ \hline
Group Discussion ALP&4&14&56&\cellcolor{yellowlearn}Experiential Learning\\ \hline
Group Discussion Social Project&4&14&56&\cellcolor{yellowlearn}Experiential Learning\\ \hline
\multirow{2}{*}{Kegiatan Khusus}&2&14&28&\cellcolor{greenlearn}Social Learning\\ \cline{2-5}
&6&14&48&\cellcolor{yellowlearn}Experiential Learning\\ \hline
Business Immersion&40&9&360&\cellcolor{yellowlearn}Experiential Learning\\ \hline
\end{tabularx}\end{center}

\begin{visualexplanation}
Tabel mempertahankan setiap baris alokasi sesi dari dokumen sumber. Beberapa kegiatan menempati dua baris karena satu nama kegiatan berkontribusi pada dua kategori: Pre-Learning terbagi menjadi formal dan social learning; Case Study terbagi menjadi social dan experiential learning; Kegiatan Khusus terbagi menjadi social dan experiential learning. Business Immersion memiliki intensitas terbesar, yaitu 40 sesi per minggu selama 9 minggu atau 360 sesi, dan seluruhnya ditempatkan pada experiential learning. Nilai 48 pada baris experiential Kegiatan Khusus dipertahankan persis seperti data sumber.
\end{visualexplanation}

\Needspace{8cm}
\subsection{Proporsi 70--20--10}
\begin{center}
\begin{tikzpicture}
\def\R{2.35}\def\r{1.28}
\fill[yellowlearn] (0,0)--(90:\R) arc (90:342:\R)--cycle;
\fill[greenlearn] (0,0)--(342:\R) arc (342:414:\R)--cycle;
\fill[cyanlearn] (0,0)--(54:\R) arc (54:90:\R)--cycle;
\fill[white] (0,0) circle (\r);
\node[font=\bfseries] at (216:1.82) {70\%};
\node[font=\bfseries] at (18:1.82) {20\%};
\node[font=\bfseries] at (72:1.82) {10\%};
\node[anchor=west,text width=7.8cm] at (3.2,1.65) {\colorbox{cyanlearn}{\strut\textit{Formal Learning: Structuring}}\\Memberikan basis dan pondasi pembelajaran teoritikal secara formal yang mendukung aplikasi praktis};
\node[anchor=west,text width=7.8cm] at (3.2,0) {\colorbox{greenlearn}{\strut\textit{Social Learning: Interacting}}\\Menghubungkan ilmu dengan \textit{feedback} melalui interaksi, diskusi kolaboratif, dan pengalaman grup bersama};
\node[anchor=west,text width=7.8cm] at (3.2,-1.75) {\colorbox{yellowlearn}{\strut\textit{Experiential Learning: Doing}}\\Transformasi ilmu menjadi \textit{skill} dengan pembelajaran melalui pengambilan keputusan dan evaluasi aksi nyata};
\end{tikzpicture}
\end{center}
\begin{visualexplanation}
Donut chart membagi keseluruhan model menjadi tiga irisan. Irisan terbesar, 70\%, adalah experiential learning atau \textit{doing}: peserta mengubah ilmu menjadi keterampilan melalui keputusan dan evaluasi aksi nyata. Irisan 20\% adalah social learning atau \textit{interacting}: pembelajaran diperkuat melalui feedback, diskusi, dan pengalaman grup. Irisan 10\% adalah formal learning atau \textit{structuring}: teori memberi basis yang diperlukan sebelum diterapkan. Warna kategori pada tabel di atas sama dengan warna irisan, sehingga setiap alokasi sesi dapat dibaca sebagai kontribusi kepada komposisi 70--20--10.
\end{visualexplanation}

\section{Learning Architecture and Methodology}
\subsection{Learning Architecture: dari Cohort ke Small Group}
\begin{center}
\begin{tikzpicture}[level/.style={rounded corners=1.5mm,minimum height=13mm,text=white,
 align=center,font=\small},>={Latex[length=3mm]},node distance=4mm]
\node[level,fill=navy,text width=24.5cm] (co) {\textcolor{yellowlearn}{\bfseries COHORT} -- \textbf{399 Peserta}\\Executive Seminar, Guest Lecture CEO BUMN, Panel Discussion, Reflection Session};
\node[level,fill=navy!92,text width=20.7cm,below=of co] (kp) {\textcolor{yellowlearn}{\bfseries KELAS PARALEL} -- \textbf{8 Kelas @40--70 Peserta}\\Business Class, Leadership Class, Case Study, Business Game, Simulation};
\node[level,fill=navy!82,text width=17.4cm,below=of kp] (sg) {\textcolor{yellowlearn}{\bfseries SMALL GROUP} -- \textbf{40 Kelompok @ $\pm$10 Peserta}\\Action Learning Project (ALP), Business Coaching 3x, Leadership Mentoring 3x, Social Project};
\node[level,fill=blue!82,text width=13.5cm,below=of sg] (ind) {\textcolor{yellowlearn}{\bfseries INDIVIDUAL} -- \textbf{399 Peserta}\\Self Learning, Assignment, Assessment per Modul, Individual Report};
\draw[->,very thick,blue!65] (co)--(kp); \draw[->,very thick,blue!65] (kp)--(sg); \draw[->,very thick,blue!65] (sg)--(ind);
\end{tikzpicture}
\end{center}
\begin{visualexplanation}
Arsitektur menyempit dari kegiatan berskala cohort ke pekerjaan individual. Seluruh 399 peserta berkumpul untuk seminar, kuliah tamu CEO BUMN, diskusi panel, dan reflection session. Mereka kemudian dibagi menjadi 8 kelas paralel berisi 40--70 peserta untuk business class, leadership class, case study, business game, dan simulation. Kelas dipecah lagi menjadi 40 small group beranggotakan sekitar 10 peserta untuk ALP, tiga business coaching, tiga leadership mentoring, dan social project. Pada lapisan terakhir, seluruh 399 peserta tetap bertanggung jawab secara individual atas self learning, assignment, assessment per modul, dan individual report. Penyempitan lebar kotak memperlihatkan berkurangnya ukuran kelompok sekaligus meningkatnya personalisasi dan akuntabilitas.
\end{visualexplanation}

\subsection{Core Learning Methodology}
\begin{tcbraster}[raster columns=2,raster equal height=rows,raster column skip=3mm]
\begin{infobox}[Learning Methodology 1: 6-Core Courses]{navy}
MK1 -- Organizational Behavior \& Managing People\par MK2 -- Financial Management \& Strategy\par
MK3 -- Marketing Management\par MK4 -- Operations \& Supply Chain Management\par
MK5 -- Decision Making \& Negotiation\par MK6 -- Business Analytics
\end{infobox}
\begin{infobox}[Learning Methodology 2: ALP + Business Immersion + Social Project]{navy}
Tantangan pada bisnis nyata dengan dampak terukur.
\end{infobox}
\end{tcbraster}

\subsection{Why Action Learning Project (ALP)?}
\begin{tcolorbox}[colback=lightgray,colframe=lightgray,arc=2mm]\centering
\textbf{Action Learning Project (ALP) -- ``Learning through real business challenges with measurable impact''}\par\medskip
\begin{tikzpicture}[value/.style={draw=navy,fill=white,rounded corners=2mm,minimum width=7.5cm,
 minimum height=12mm,align=center,font=\bfseries},>={Latex[length=2.5mm]}]
\node[value] (v1) at (0,0) {Equip \& Transform Participants' Capabilities};
\node[value] (v2) at (8.5,0) {Real Business ROI that Creates Value};
\node[value] (v3) at (17,0) {True Test of Agility \& Resilience};
\draw[->,thick,navy] (v1)--(v2); \draw[->,thick,navy] (v2)--(v3);
\end{tikzpicture}
\end{tcolorbox}
\begin{tcbraster}[raster columns=2,raster equal height=rows,raster column skip=3mm]
\begin{infobox}[The Carpenter Analogy]{navy}
\textit{``Tukang bangunan yang tahu menggunakan alat, belum tentu bisa membuat rumah. Hanya pengalaman dalam membuat rumah maka ia dapat menggunakan keahliannya.''}
\begin{itemize}
\item Can this person manage ambiguity, not just structured tasks?
\item Can they work together under real-time pressure?
\item Are they ready for a bigger, less-defined mandate?
\end{itemize}
\end{infobox}
\begin{infobox}[An ALP is a ``stretch'' assignment]{navy}
\begin{itemize}[label=$\checkmark$]
\item A clear picture of the problem, grounded in data and interviews
\item A named set of options, with trade-offs made explicit
\item A prioritized, financially-aware recommendation
\item A view on what it takes to implement it
\item Producing analysis and concrete, actionable recommendations a sponsor can act on
\end{itemize}
\end{infobox}
\end{tcbraster}
\begin{visualexplanation}
ALP connects the two learning methodologies. The six core courses supply tools and concepts, while ALP, Business Immersion, and the Social Project apply them to a measurable business challenge. The three linked value boxes show the intended progression: participant capability is transformed through practice; that practice must create real business return; and the assignment tests agility and resilience under ambiguity. The carpenter analogy explains why knowing individual tools is insufficient without building something complete. The stretch-assignment panel then defines the expected evidence: a data-grounded problem, explicit options and trade-offs, a financially aware recommendation, an implementation view, and concrete actions a sponsor can use.
\end{visualexplanation}

\section{Learning Journey Timeline: 6-Core Courses}
\textit{6 Mata Kuliah dirancang untuk membekali pelajar dengan kemampuan Business Leadership, dilaksanakan di Kampus Nawasena.}
\begin{tcbraster}[raster columns=3,raster equal height=rows,raster column skip=2mm,raster row skip=2mm]
\begin{infobox}[MK1 -- Organizational Behavior and Managing People]{blue}14 Lecture @ 1 Sesi\qquad 14 Case @ 2 Sesi\end{infobox}
\begin{infobox}[MK2 -- Financial Management and Strategy]{teal}14 Lecture @ 1 Sesi\qquad 14 Case @ 2 Sesi\end{infobox}
\begin{infobox}[MK3 -- Marketing Management]{maroon}14 Lecture @ 1 Sesi\qquad 13 Case @ 2 Sesi\qquad 1 Game @ 2 Sesi\end{infobox}
\begin{infobox}[MK4 -- Operations and Supply Chain Management]{brick}14 Lecture @ 1 Sesi\qquad 14 Case @ 2 Sesi\qquad 1 Game @ 2 Sesi\end{infobox}
\begin{infobox}[MK5 -- Decision Making and Negotiation]{orange}14 Lecture @ 1 Sesi\qquad 7 Case @ 2 Sesi\qquad 7 Roleplay @ 2 Sesi\end{infobox}
\begin{infobox}[MK6 -- Business Analytics]{gold}14 Lecture @ 1 Sesi\qquad 14 Case @ 2 Sesi\end{infobox}
\end{tcbraster}
\subsection{Alur Mingguan per MK}
\begin{center}\scriptsize\setlength{\tabcolsep}{2pt}
\begin{tabular}{|L{3.5cm}|*{14}{M{1.35cm}|}}\hline
\textbf{Alur Mingguan per MK}&\textbf{M1}&\textbf{M2}&\textbf{M3}&\textbf{M4}&\textbf{M5}&\textbf{M6}&\textbf{M7}&\textbf{M8}&\textbf{M9}&\textbf{M10}&\textbf{M11}&\textbf{M12}&\textbf{M13}&\textbf{M14}\\ \hline
\textbf{Perkuliahan}&\multicolumn{7}{c|}{\cellcolor{blue!78}\color{white}Blok 1 -- lecture tentang topik pembelajaran @ 1 sesi}&\multicolumn{7}{c|}{\cellcolor{blue!78}\color{white}Blok 2 -- lecture tentang topik pembelajaran @ 1 sesi}\\ \hline
\textbf{Case Study}&\multicolumn{7}{c|}{\cellcolor{blue}\color{white}Intensif diskusi kasus @ 2 sesi}&\multicolumn{7}{c|}{\cellcolor{blue}\color{white}Intensif diskusi kasus @ 2 sesi}\\ \hline
\textbf{Gamification}&&&&&&&\cellcolor{red}\color{white}\shortstack{OSCM\\2 Sesi}&&&&&&&\cellcolor{red}\color{white}\shortstack{Mkt\\2 Sesi}\\ \hline
\textbf{Ujian}&&&&&&&&\cellcolor{orange!82}\color{white}\shortstack{Mid\\Term}&&&&&&\cellcolor{orange!82}\color{white}\shortstack{Final\\Term}\\ \hline
\textbf{Guest Lecture CEO BUMN}&&&\cellcolor{midgray}\color{white}1&&&&&\cellcolor{midgray}\color{white}2&&&&&&\cellcolor{midgray}\color{white}3\\ \hline
\rowcolor{lightblue}\textbf{Titik Sinkronisasi}&\shortstack{\tiny\textcolor{navy}{$\bullet$}\\\tiny Kick Off\\\tiny ALP}&&&&\shortstack{\tiny\textcolor{navy}{$\bullet$}\\\tiny Milestone 1}&&&&&&\shortstack{\tiny\textcolor{navy}{$\bullet$}\\\tiny Milestone 2}&&&\shortstack{\tiny\textcolor{navy}{$\bullet$}\\\tiny Final\\\tiny Presentation}\\ \hline
\end{tabular}\end{center}
\begin{visualexplanation}
Enam kotak mata kuliah merangkum jumlah lecture, case, game, atau roleplay yang dimiliki setiap MK. Timeline menjelaskan pola 14 minggu: lecture dan case study berjalan dalam dua blok tujuh-minggu; gamification OSCM berada pada M7 dan gamification Marketing pada M14; ujian mid term berada pada M8 dan final term pada M14; Guest Lecture CEO BUMN muncul pada M3, M8, dan M14. Baris titik sinkronisasi menghubungkan kegiatan mata kuliah dengan ALP melalui Kick Off pada M1, Milestone 1 pada M5, Milestone 2 pada M11, dan Final Presentation pada M14.
\end{visualexplanation}

\Needspace{11cm}
\section{Learning Journey Timeline: Action Learning Project}
\textit{Enam tahap ALP SMART-gated di 10 BUMN mitra, dengan tiga milestone sebagai gerbang antar tahap ALP.}
\begin{tcbraster}[raster columns=3,raster equal height=rows,raster column skip=2mm,raster row skip=2mm]
\begin{infobox}[1. Problem Scoping \hfill M2--3]{green}Identifikasi \textit{business issue} bersama \textit{project owner}.\end{infobox}
\begin{infobox}[2. Data Gathering \hfill M4--6]{green}Interview \textit{stakeholder} dan survei \textit{customer}.\end{infobox}
\begin{infobox}[3. Analysis \hfill M7--9]{green}Integrasi strategi, finansial, SDM, operasi.\end{infobox}
\begin{infobox}[4. Solution Design \hfill M10--12]{green}Rekomendasi implementatif dan terukur.\end{infobox}
\begin{infobox}[5. Coaching Review \hfill M13--14]{green}Review bersama \textit{coach} dan \textit{project owner}.\end{infobox}
\begin{infobox}[6. Final Presentation \hfill M15--16]{green}Presentasi ke manajemen dan panel.\end{infobox}
\end{tcbraster}
\subsection{Alur Mingguan}
\begin{center}\scriptsize\setlength{\tabcolsep}{1.7pt}
\begin{tabular}{|L{3.8cm}|*{16}{M{1.15cm}|}}\hline
\textbf{Alur Mingguan}&\textbf{M1}&\textbf{M2}&\textbf{M3}&\textbf{M4}&\textbf{M5}&\textbf{M6}&\textbf{M7}&\textbf{M8}&\textbf{M9}&\textbf{M10}&\textbf{M11}&\textbf{M12}&\textbf{M13}&\textbf{M14}&\textbf{M15}&\textbf{M16}\\ \hline
\textbf{Tahap ALP}&&\multicolumn{2}{c|}{\cellcolor{greenlearn}Scoping}&\multicolumn{3}{c|}{\cellcolor{greenlearn}Data Gathering}&\multicolumn{3}{c|}{\cellcolor{greenlearn}Analysis}&\multicolumn{3}{c|}{\cellcolor{greenlearn}Solution Design}&\multicolumn{2}{c|}{\cellcolor{greenlearn}Coaching Review}&\multicolumn{2}{c|}{\cellcolor{greenlearn}Final Presentation}\\ \hline
\textbf{Mentoring Praktisi BUMN}&&&\cellcolor{blue}\color{white}1&&&&\cellcolor{blue}\color{white}2&&&&\cellcolor{blue}\color{white}3&&&&&\\ \hline
\textbf{Coaching Akademisi}&&&&\cellcolor{blue}\color{white}1&&&&\cellcolor{blue}\color{white}2&&&&\cellcolor{blue}\color{white}3&&&&\\ \hline
\textbf{Milestone}&&&&&&\cellcolor{orange!82}\color{white}M1&&&&&&\cellcolor{orange!82}\color{white}M2&&&&\cellcolor{orange!82}\color{white}M3\\ \hline
\textbf{Immersion \& Social Project}&&&&&\cellcolor{midgray}\color{white}Imm 1&&&&&&\cellcolor{midgray}\color{white}Imm 2&&\multicolumn{3}{c|}{\cellcolor{midgray}\color{white}Social Project}&\\ \hline
\rowcolor{lightblue}\textbf{Titik Sinkronisasi}&&\shortstack{\tiny\textcolor{navy}{$\bullet$}\\\tiny Kick Off\\\tiny ALP}&&&&\shortstack{\tiny\textcolor{navy}{$\bullet$}\\\tiny Milestone 1}&&\shortstack{\tiny\textcolor{navy}{$\bullet$}\\\tiny Mid Term}&&&&\shortstack{\tiny\textcolor{navy}{$\bullet$}\\\tiny Milestone 2}&&&&\shortstack{\tiny\textcolor{navy}{$\bullet$}\\\tiny Final\\\tiny Presentation}\\ \hline
\end{tabular}\end{center}
\begin{visualexplanation}
Enam kartu tahap memberi arti pada blok waktu dalam tabel. Scoping berlangsung pada M2--3, Data Gathering pada M4--6, Analysis pada M7--9, Solution Design pada M10--12, Coaching Review pada M13--14, dan Final Presentation pada M15--16. Mentoring Praktisi BUMN ditempatkan pada M3, M7, dan M11; Coaching Akademisi pada M4, M8, dan M12. Milestone M1, M2, dan M3 berada pada M6, M12, dan M16. Immersion 1 berlangsung pada M5, Immersion 2 pada M11, dan Social Project pada M13--15. Baris sinkronisasi menandai Kick Off ALP pada M2, Milestone 1 pada M6, Mid Term pada M8, Milestone 2 pada M12, dan Final Presentation pada M16.
\end{visualexplanation}